% Applied Software Development Lab -- Lab Handout 11: Levels, lives, variety, and control
% Copyright (c) Michael Lehn. Licensed under CC BY-SA 4.0.
%
% Region coordinates (bp from the top of the page) follow the page layout of
% the original German handout "Laborblatt 11".
\documentclass{aslab}
\sethandout{11}{Levels, lives, variety, and control}
\graphicspath{{../figures/handout11/}{../figures/shared/}}

%% Python code blocks: one source line per output line, spaces are kept
%% (each one character wide) and the baseline distance is fixed.
%% Syntax colors: \kw keyword, \fn function/builtin/constant, \num number or
%% string, \str string (dark red), \cmt comment, \imp import statement.
\definecolor{codestring}{HTML}{B51700}
\newcommand*\kw[1]{\textcolor{codeorange}{#1}}
\newcommand*\fn[1]{\textcolor{codecyan}{#1}}
\newcommand*\num[1]{\textcolor{codered}{#1}}
\newcommand*\str[1]{\textcolor{codestring}{#1}}
\newcommand*\cmt[1]{\textcolor{codepurple}{#1}}
\newcommand*\imp[1]{\textcolor{codemagenta}{#1}}
\makeatletter
\newcommand*\py@space{\leavevmode\hbox to 0.6em{}}
{\catcode`\^^M=\active\catcode`\ =\active%
 \gdef\py@setup{\def^^M{\leavevmode\par}\let =\py@space}%
 \gdef\py@gobble#1^^M{}}
\newenvironment{pycode}{\par\ttfamily\hyphenchar\font=-1
  \lineskiplimit=-\maxdimen\lineskip=0pt\parindent=0pt
  \parfillskip=0pt plus 1fil\relax
  \catcode`\^^M=\active\catcode`\ =\active\py@setup\py@gobble}{}
\makeatother
%% Emoji at the start of a line, followed by the text at +18bp
\newcommand*\emo[1]{\makebox[18bp][l]{\emoji{#1}}}
\newcommand*\catemo{\makebox[18bp][l]{{\emojifont ^^^^^^01f408^^^^200d^^^^2b1b}}}

\begin{document}
%% ===================================================================== 1 ==
\maketitleblock

\begin{region}[pitch=15.5]{118.0}
{\bfseries\color{aslabblue}Learning objective: Bringing a software project to a
playable, impressive conclusion under realistic conditions (read: shortly
before the deadline)\par}
\end{region}

\begin{region}[pitch=14.27]{154.1}
{\bfseries\color{aslabblue}Note: This time, the steps contain no explicit
prompts to commit your current state to the repository. You can now handle
that on your own -- and rightly so!\\
In case you do hesitate for a moment: Step~13 on the last page has a little
reminder and instructions.\par}
\end{region}

\begin{region}{225.4}
\heading{Preparation}
\end{region}

\begin{region}{243.1}
As usual, we work with two terminals:
\end{region}

\begin{region}[pitch=14.3]{256.2}
\begin{dashlist}
\item \textbf{left}: edit the source code
\item \textbf{right}: test the game
\end{dashlist}
\end{region}

\begin{region}{290.3}
Cloning the repository:
\end{region}

\begin{region}[pitch=13.2]{309.2}
\begin{pycode}
git clone https://gitlab.uni-ulm.de/USERNAME/asp-ulm
cd asp-ulm/python
python breakout.py
\end{pycode}
\end{region}

\begin{region}{355.0}
As always, \textbf{take a moment to familiarize yourselves} with the features
your game currently supports.
\end{region}

\begin{region}[pitch=14.2]{374.8}
If \textbf{unexpected crashes occur} -- and you are \textbf{lucky enough} to
discover them now -- it is best to fix them right away. \textbf{We are happy
to help.}
\end{region}

\begin{region}{416.5}
\heading{Step 1: Bombs are not just for decoration}
\end{region}

\begin{region}[pitch=13.25]{435.1}
When the paddle is hit by a bomb, it should be \textbf{out of action} for a
certain time. In this step, we first only check whether a bomb hits the
paddle. If so, the bomb should detonate.
\end{region}

\begin{region}[pitch=13.9]{480.1}
To do this, you can \textbf{proceed similarly} to the check whether a ball
hits the paddle.\\
If you want to do it really well: take a look at \code{ball.py}. There, the
\code{move} method checks whether the \textbf{center} of the ball crosses the
\textbf{top edge} of the paddle.
\end{region}

\begin{region}[pitch=13.75]{525.9}
If you proceed analogously, you can adapt the \code{move} method in
\code{bomb.py} accordingly. The code is even \textbf{a bit simpler} here,
because our bombs always fall \textbf{straight down}.
\end{region}

\begin{region}[pitch=14.2]{571.2}
When the paddle is hit, you simply set the attribute \code{self.hit = True}.
The \code{draw} method takes care of the rest -- \textbf{without any
changes}.
\end{region}

\begin{region}{603.4}
Your \code{move} method in \code{bomb.py} could then look like this, for
example:
\end{region}

\begin{region}[]{630.0}
\begin{pycode}
    \kw{def} \fn{move}(self):
        \kw{if} self.y < \num{0} \kw{or} self.y > Game.height:
            self.alive = \fn{False}
            \kw{return}
        new\_y = self.y + self.vy
        \kw{if} self.y < Paddle.y-\num{10} \kw{and} new\_y >= Paddle.y-\num{10}:
            \kw{if} self.x >= Paddle.x \kw{and} self.x <= Paddle.x + Paddle.w:
                self.hit = \fn{True}
        self.y = new\_y
\end{pycode}
\end{region}

\begin{region}{752.7}
\textbf{Task:} Test whether, after your change, the bombs explode when they
hit the paddle.
\end{region}

\nextpage
%% ===================================================================== 2 ==
\begin{region}{69.4}
\heading{Step 2: Bombs are bad for the paddle}
\end{region}

\begin{region}[pitch=14.0]{88.1}
When the paddle has been hit, this should have \textbf{consequences}: for a
certain time, it is \textbf{out of action} -- and, in particular,
\textbf{not visible}. Here is a small \textbf{template} showing how the
feature can be implemented \textbf{at least partially}. You will extend the
rest yourselves in the course of the task.
\end{region}

\begin{region}{151.9}
\emo{1F9F1}\textbf{In} \code{paddle.py}
\end{region}

\begin{region}[pitch=14.2]{167.9}
The paddle gets a new static attribute \code{damage}, which is initialized
with \code{0}. It indicates \textbf{for how many frames} the paddle cannot be
used. The class \code{Paddle} thus begins like this, for example:
\end{region}

\begin{region}[]{200.1}
\begin{pycode}
\kw{class} \fn{Paddle}:
    c = (\num{0}, \num{0}, \num{255})
    w, h = \num{100}, \num{30}
    x, y = \fn{None}, \fn{None}
    damage = \num{0}
\end{pycode}
\end{region}

\begin{region}{268.6}
\emo{1F4A5}\textbf{New method} \code{hit}
\end{region}

\begin{region}[pitch=13.3]{285.3}
The method \code{hit()} ``destroys'' the paddle -- \textbf{but only if it is
still intact}. If it is already damaged, nothing happens.
\end{region}

\begin{region}[pitch=13.2]{317.5}
\begin{pycode}
    \kw{def} \fn{hit}():
        \kw{if} Paddle.damage == \num{0}:
\end{pycode}
\end{region}
\begin{region}[]{342.7}
\begin{pycode}
            Paddle.damage = \num{90}
            \kw{return} \fn{True}
        \kw{return} \fn{False}
\end{pycode}
\end{region}

\begin{region}[pitch=13.6]{383.4}
\textbf{Note:} The method returns \code{True} if the paddle has just been
damaged. Otherwise, it returns \code{False} (e.g., if it is already broken
anyway).
\end{region}

\begin{region}{414.6}
\emo{270F}\textbf{Adapting the} \code{draw} \textbf{method}
\end{region}

\begin{region}{430.4}
If the paddle is damaged, it is not drawn -- but the damage is reduced:
\end{region}

\begin{region}[]{443.8}
\begin{pycode}
    \kw{def} \fn{draw}():
        \kw{if} Paddle.damage > \num{0}:
            Paddle.damage -= \num{1}
            \kw{return}
        pygame.draw.rect(Game.screen, Paddle.c,
                         (Paddle.x, Paddle.y, Paddle.w, Paddle.h))
\end{pycode}
\end{region}

\begin{region}{516.7}
\subheading{Task}
\end{region}

\begin{region}[pitch=13.38]{528.8}
\begin{tasks}
\item Adapt the code in \code{paddle.py} as described above. As always, try
  to understand exactly what happens. Now, of course, \code{Paddle.hit()}
  should be called when an intact paddle is hit by a bomb.\\
  To do this, change the line \code{self.hit = True} in \code{bomb.py} to
  \code{self.hit = Paddle.hit()}. This way, \code{self.hit} is only set to
  \code{True} if the paddle was actually damaged by the bomb hit.
\end{tasks}
\end{region}

\begin{region}[pitch=13.2]{608.9}
\begin{tasks}[start=2]
\item The paddle is now invisible for a while after it has been hit -- but
  it is still there. That is, even in its damaged state, it can still
  reflect balls.\\
  Hence, two subtasks:
\end{tasks}
\end{region}

\begin{region}[x=74.7,right=0,pitch=13.95]{649.4}
\begin{steps}[label=\textbf{(\alph*)}]
\item \textbf{(hard) Observe the behavior}\\
  First, try to \textbf{consciously observe} the problem before you fix it.
  \textbf{Can you make a ball bounce off an invisible paddle?} That is not
  easy at all -- but this is exactly the kind of behavior a good test team
  should discover
\end{steps}
\end{region}

\begin{region}[x=74.7,right=0,pitch=14.0]{705.3}
\begin{steps}[label=\textbf{(\alph*)},start=2]
\item \textbf{(easy) Fix the problem}\\
  In \code{ball.py}, adapt the place where it is checked whether the ball
  collides with the paddle. Add an additional condition there so that the
  ball \textbf{is only reflected} if the paddle \textbf{is intact}.
\end{steps}
\end{region}

\begin{region}[pitch=13.2]{760.7}
\begin{tasks}[start=3]
\item Now also fix the problem that you can currently still fire rockets
  with a destroyed paddle.
\end{tasks}
\end{region}

\nextpage
%% ===================================================================== 3 ==
\begin{region}{69.4}
\heading{Step 3: Long live variety (part 1 of 3)}
\end{region}

\begin{region}[pitch=14.3]{88.1}
Currently, our bricks are still rather \textbf{monotonous}: they are either
\textbf{there or gone} -- nothing more.\\
That is about to change. There should be \textbf{different kinds of bricks},
for example:
\end{region}

\begin{region}[pitch=14.3]{115.6}
\begin{dashlist}
\item Bricks that you have to \textbf{hit twice} (only then do they
  disappear)
\item Bricks that bring \textbf{a new ball} into the game when destroyed
\item Bricks that give the paddle \textbf{a temporary protective shield} --
  after that, you are immune to bombs for a while
\end{dashlist}
\end{region}

\begin{region}{171.1}
\begin{dashlist}
\item And maybe more\dots
\end{dashlist}
\end{region}

\begin{region}{199.0}
\subheading{Idea: numbers for brick types}
\end{region}

\begin{region}{216.7}
The type of a brick can initially be encoded with numbers:
\end{region}

\begin{region}[pitch=14.67]{230.0}
\begin{dashlist}
\item \code{0} $\rightarrow$ no brick
\item \code{1} $\rightarrow$ normal brick (what we have had so far)
\item \code{2} $\rightarrow$ brick that has to be hit twice
\item \code{3} $\rightarrow$ brick that creates a new ball
\item \code{4} $\rightarrow$ brick that activates a protective shield
\end{dashlist}
\end{region}

\begin{region}[pitch=13.3]{316.7}
This has the advantage that, as before, you can work with a \code{numpy}
array or with comparison operators.
\end{region}

\begin{region}[pitch=13.3]{357.3}
\textbf{But}: \textbf{numbers alone} are \textbf{not particularly
expressive} -- you quickly forget what 3 or 4 stood for.
\end{region}

\begin{region}{397.9}
\textbf{Remedy:} \code{enum} \textbf{types}
\end{region}

\begin{region}{416.5}
Here comes a mini crash course -- entirely in the style of:
\textbf{``Learning by using and playing around.''}
\end{region}

\begin{region}{443.3}
Create a new file named \code{bricktype.py} with the following content:
\end{region}

\begin{region}[pitch=13.2]{470.0}
\begin{pycode}
  \kw{1} \imp{from} enum \imp{import} IntEnum
  \kw{2}
\end{pycode}
\end{region}
\begin{region}[]{495.2}
\begin{pycode}
  \kw{3 class} \fn{BrickType}(IntEnum):
  \kw{4}     NONE = \num{0}
  \kw{5}     NORMAL = \num{1}
  \kw{6}     BALL = \num{2}
  \kw{7}     SHIELD = \num{3}
  \kw{8}     TWICE = \num{4}
  \kw{9}
 \kw{10}     \kw{def} \fn{label}(self):
 \kw{11}         \kw{return} f\num{"Dummy label for \{self\}"}
 \kw{12}
 \kw{13 def} \fn{main}():
 \kw{14}     \fn{print}(\fn{max}(BrickType))
 \kw{15}     \fn{print}(BrickType.BALL)
 \kw{16}     \fn{print}(BrickType.BALL.label())
 \kw{17}     foo = BrickType(\num{4})
 \kw{18}     \fn{print}(foo)
 \kw{19}     \fn{print}(foo.label())
 \kw{20}
 \kw{21}
 \kw{22 if} \_\_name\_\_ == \num{"\_\_main\_\_"}:
 \kw{23}     main()
\end{pycode}
\end{region}

\nextpage
%% ===================================================================== 4 ==
\begin{region}{69.4}
\subheading{Explanations}
\end{region}

\begin{region}[pitch=14.25]{95.9}
\begin{dashlist}
\item In \textbf{line 3}, \code{BrickType} is declared as an extension of
  \code{IntEnum}. At this point, we will \textbf{not go into detail} about
  what this means -- \textbf{except} in a later \textbf{subtask}, in which you
  remove the \code{(IntEnum)} once to see \textbf{what suddenly no longer
  works}.
\item In \textbf{lines 4--8}, the different \textbf{kinds of bricks} are
  defined as constants with \textbf{meaningful names} -- much more readable
  than plain numbers.
\item In \textbf{lines 10--11}, a method label is defined that, for now, only
  returns a \textbf{placeholder text} -- but already shows \textbf{what is
  possible with enums}.
\end{dashlist}
\end{region}

\begin{region}{222.0}
\subheading{Task}
\end{region}

\begin{region}[pitch=13.5]{248.8}
\begin{tasks}
\item Try it out with \code{python bricktype.py}. Observe the following:
  \begin{bullets}
  \item \code{max(BrickType)} returns the largest enum member -- in this
    case, \code{BrickType.TWICE}
  \item You can use \code{BrickType.BALL} instead of the number \code{2} --
    much more expressive!
  \item \code{BrickType.BALL} is not simply a number: for example, you can
    apply \code{.label()} to it
  \item With \code{BrickType(4)}, you can convert an ordinary integer value
    back into a \code{BrickType} object
  \end{bullets}
\end{tasks}
\end{region}

\begin{region}[pitch=14.3]{343.3}
\begin{tasks}[start=2]
\item \textbf{What do we gain from this} \code{(IntEnum)}\textbf{?}\\
  In \textbf{line 3}, remove the \code{(IntEnum)} and run the program again:
  \begin{bullets}
  \item \textbf{Error in line 14}. The function \code{max()} can no longer be
    applied to the class.
  \end{bullets}
\end{tasks}
\end{region}

\begin{region}[x=74.7,right=0,pitch=13.6]{391.7}
\textbf{Comment out line 14} and run the program again:
\begin{bullets}
\item Now you get a new \textbf{error in line 16}. \code{BrickType.BALL} is
  now just a number, and the method \code{label} can no longer be applied.
\end{bullets}
\end{region}

\begin{region}[x=74.7,right=0,pitch=13.5]{439.2}
\textbf{What we learn about the effect of} \code{(IntEnum)}\textbf{:}
\begin{bullets}
\item Attributes such as \code{BrickType.BALL} are not ordinary \code{int}
  values but objects of type \code{BrickType}
\item These objects have methods such as \code{.label()}, which you define
  yourself
\item In addition, you ``inherit'' certain things, e.g., that you can write
  \code{max(BrickType)}
\end{bullets}
\end{region}

\begin{region}[x=74.7,right=0,pitch=14.2]{526.2}
\textbf{Undo all changes}: you are welcome to look into the exact details on
a \textbf{quieter day} -- today is about \textbf{understanding by trying
things out}, not about the last Python detail \emoji{1F609}
\end{region}

\begin{region}[pitch=13.3]{574.8}
\begin{tasks}[start=3]
\item \textbf{Extending the} \code{main()} \textbf{function and improving}
  \code{label()}\\
  Change the function main() as follows:
\end{tasks}
\end{region}

\begin{region}[x=74.7,pitch=13.2]{607.1}
\begin{pycode}
\kw{def} \fn{main}():
    \kw{for} i \kw{in} \fn{range}(\num{0}, \fn{max}(BrickType)):
        bt = BrickType(i)
        \fn{print}(f\str{"i = \{i\}, bt = \{bt\}, bt.label = \{bt.label()\}"})
\end{pycode}
\end{region}

\begin{region}[x=74.7,right=0,pitch=13.6]{665.3}
Run the program and follow what happens. Now change the method
\code{label()} so that -- depending on the value of \code{self} -- it
returns the following text:
\end{region}

\begin{region}[x=74.7,right=0,pitch=13.5]{692.4}
\begin{bullets}
\item \code{"Ball"} for \code{BrickType.BALL}
\item \code{"Shield"} for \code{BrickType.SHIELD}
\item \code{"Twice"} for \code{BrickType.TWICE}
\item In all other cases, the empty string \code{""}
\end{bullets}
\end{region}

\begin{region}[x=74.7,right=0]{750.5}
\emo{1F449}This is a small exercise in \code{if}--\code{elif}--\code{else}.
\end{region}

\nextpage
%% ===================================================================== 5 ==
\begin{region}{69.4}
\heading{Step 4: Long live variety (part 2 of 3)}
\end{region}

\begin{region}[pitch=13.8]{88.1}
Now the class BrickType is \textbf{brought into the game very gently}. At
first, it is \textbf{only about the display}: bricks should \textbf{differ
visually} by having a small \textbf{label} text printed on them.
\end{region}

\begin{region}{129.8}
\textbf{Important}: these bricks \textbf{do not behave differently yet} --
that comes later.
\end{region}

\begin{region}{160.7}
\emo{1F9F1}In \code{wall.py}
\end{region}

\begin{region}{189.8}
First, import our new class with \code{from bricktype import BrickType}.
\end{region}

\begin{region}{220.6}
\emo{1F6E0}\textbf{Adapting the} \code{init} \textbf{method}
\end{region}

\begin{region}{236.6}
Then change the method \code{init} as follows:
\end{region}

\begin{region}[]{263.3}
\begin{pycode}
 \kw{12}     \kw{def} \fn{init}():
 \kw{13}         \kw{for} i \kw{in} \fn{range}(\num{4}, \num{8}):
 \kw{14}             \kw{for} j \kw{in} \fn{range}(\num{0}, Game.brick\_cols):
 \kw{15}                 Wall.brick[i, j] = random.randint(\num{1}, \fn{max}(BrickType))
\end{pycode}
\end{region}

\begin{region}[pitch=13.5]{324.5}
Previously, all cells there were bluntly filled with \code{1} -- now each
cell gets a random \code{BrickType} value (from \code{1} up to and including
\code{4}).
\end{region}

\begin{region}{368.8}
\emo{1F5BC}\textbf{Extending the} \code{draw} \textbf{method}
\end{region}

\begin{region}[pitch=14.2]{384.8}
Now we make the new brick types visible. In the \code{draw} method, add the
following code starting at \textbf{line 45}:
\end{region}

\begin{region}[]{425.7}
\begin{pycode}
\kw{33}     \kw{def} \fn{draw}():
\kw{34}         \kw{for} i \kw{in} \fn{range}(\num{0}, Wall.brick.shape[\num{0}]):
\kw{35}             \kw{for} j \kw{in} \fn{range}(\num{0}, Wall.brick.shape[\num{1}]):
\kw{36}                 \kw{if} \kw{not} Wall.brick[i, j]:
\kw{37}                     \kw{continue}
\kw{38}                 w = Game.brick\_width
\kw{39}                 h = Game.brick\_height
\kw{40}                 x = j * w
\kw{41}                 y = i * h
\kw{42}                 r = (x + \num{1}, y + \num{1}, w - \num{1}, h - \num{1})
\kw{43}                 pygame.draw.rect(Game.screen, (\num{223}, \num{105}, \num{165}), r)
\kw{44}
\kw{45}                 label = BrickType.label(Wall.brick[i, j])
\kw{46}                 font = pygame.font.SysFont(\fn{None}, \num{24})
\kw{47}                 text = font.render(label, \fn{True}, (\num{255}, \num{255}, \num{255}))
\kw{48}                 text\_rect = text.get\_rect(center=(x+w//\num{2}, y+h //\num{2}))
\kw{49}                 Game.screen.blit(text, text\_rect)
\end{pycode}
\end{region}

\begin{region}{656.6}
Here, the label is now drawn onto the rectangle:
\end{region}

\begin{region}[pitch=14.27]{669.8}
\begin{dashlist}
\item \textbf{Line 45}: calls \code{BrickType.label(...)} and returns an
  ordinary string.
\item \textbf{Line 46}: creates a \code{Font} object (font size 24, default
  font).
\item \textbf{Line 47}: converts the string into a \textbf{bitmap image}
  (font color: white, with antialiasing).
\item \textbf{Lines 48--49}: draws the text \textbf{centered} in the brick's
  rectangle.
\end{dashlist}
\end{region}

\nextpage
%% ===================================================================== 6 ==
\begin{region}{69.4}
\subheading{Task}
\end{region}

\begin{region}[pitch=13.5]{95.9}
\begin{tasks}
\item Test whether your adaptation works -- i.e., whether some bricks now
  have labels. You should be able to read, e.g., \code{"Ball"},
  \code{"Shield"}, or \code{"Twice"} on individual bricks.
\end{tasks}
\end{region}

\begin{region}[pitch=14.25]{123.6}
\begin{tasks}[start=2]
\item Currently, \textbf{too many special bricks} appear -- this looks a bit
  cluttered. It would be better if, with a certain probability (e.g.,
  \textbf{20\,\%}), a brick got a special function and otherwise simply
  became a \textbf{normal} brick (\code{NORMAL}).
\end{tasks}
\end{region}

\begin{region}[x=74.7]{169.8}
Suggestion:
\end{region}

\begin{region}[x=74.7,pitch=13.2]{187.6}
\begin{pycode}
\kw{if} random.randint(\num{1}, \num{100}) <= \num{20}:
    ty = random.randint(\num{1}, \fn{max}(BrickType))
\kw{else}:
    ty = \num{1}
Wall.brick[i, j] = ty
\end{pycode}
\end{region}

\begin{region}[x=74.7,right=0,pitch=13.3]{258.0}
Note: In the special selection, \code{BALL}, \code{SHIELD}, and
\code{TWICE} are then equally likely -- which is perfectly fine for our
purposes.
\end{region}

\begin{region}[pitch=13.3]{284.5}
\begin{tasks}[start=3]
\item[3.] It is not dramatic -- but still a small break in style -- that a
  new font object is created in every single \code{draw()} iteration. Even
  if the performance loss is barely measurable: good code has manners.
\end{tasks}
\end{region}

\begin{region}[x=74.7,right=0]{328.7}
It would be better if the font object were created once, centrally, in the
Game class.
\end{region}

\begin{region}[x=74.7]{352.8}
\emo{1F527}\textbf{Here's how}
\end{region}

\begin{region}[x=74.7]{373.2}
In \code{game.py}, add a new attribute to the class Game:
\end{region}

\begin{region}[x=74.7,pitch=13.2]{391.1}
\begin{pycode}
\kw{class} \fn{Game}:
    \cmt{\# ... as before ...}
    font = \fn{None}
\end{pycode}
\end{region}

\begin{region}[x=74.7]{435.4}
Then adapt the \code{init()} method as follows:
\end{region}

\begin{region}[x=74.7,pitch=13.2]{462.1}
\begin{pycode}
    \kw{def} \fn{init}():
        \cmt{\# ...}
        pygame.init()
        \cmt{\# ...}
        Game.font = pygame.font.SysFont(\fn{None}, \num{24})
\end{pycode}
\end{region}

\begin{region}[x=74.7,right=0]{542.0}
\textbf{Important:} The font object must not be created until after
\code{pygame.init()} has been called!
\end{region}

\begin{region}[x=74.7]{568.7}
Afterwards, in \code{wall.py}, you can replace the line
\end{region}

\begin{region}[x=90.0]{587.7}
\begin{pycode}
font = pygame.font.SysFont(\fn{None}, \num{24})
\end{pycode}
\end{region}

\begin{region}[x=74.7]{606.5}
with:
\end{region}

\begin{region}[x=90.0]{625.4}
\begin{pycode}
font = Game.font
\end{pycode}
\end{region}

\begin{region}[pitch=13.5]{652.8}
\begin{tasks}[start=4]
\item \textbf{Optional polish: color by type}\\
  In \code{bricktype.py}: add a method color to the class \code{BrickType}
  that -- analogous to the method \code{label()} -- returns a suitable color
  for each brick type.
\end{tasks}
\end{region}

\begin{region}[x=74.7,right=0,pitch=13.5]{701.0}
In \code{wall.py}: adapt the \code{draw()} method so that the color for the
brick is now obtained from the \code{color()} method before drawing:
\end{region}

\begin{region}[x=87.9,pitch=13.5]{741.2}
\begin{pycode}
c = BrickType.color(Wall.brick[i, j])
pygame.draw.rect(Game.screen, c, r)
\end{pycode}
\end{region}

\nextpage
%% ===================================================================== 7 ==
\begin{region}{69.4}
\heading{Step 5: Long live variety (part 3 of 3)}
\end{region}

\begin{region}[pitch=14.3]{88.1}
Our special bricks should now not \textbf{only look different} but also
\textbf{have an effect} when you hit them.
\end{region}

\begin{region}[pitch=14.3]{121.0}
We approach this \textbf{step by step} -- and start right away with the
``most complicated'' case: the \textbf{bricks that have to be hit twice}.
\end{region}

\begin{region}{158.5}
\emo{1F9E0}\textbf{Idea: changing the kind of brick on a hit}
\end{region}

\begin{region}[pitch=14.3]{175.2}
Bricks of type \code{TWICE} should \textbf{turn into} \code{NORMAL}
\textbf{on the first hit},\\
and then \textbf{disappear} on the second hit.
\end{region}

\begin{region}{212.6}
\emo{1F9F1}\textbf{In} \code{bricktype.py}
\end{region}

\begin{region}[pitch=13.3]{228.6}
Extend the class \code{BrickType} with a method \code{effect()} that returns
what a brick should turn into when it is hit:
\end{region}

\begin{region}[]{259.7}
\begin{pycode}
    \kw{def} \fn{effect}(self):
        \kw{if} self == BrickType.TWICE:
            \kw{return} BrickType.NORMAL
        \kw{else}:
            \kw{return} BrickType.NONE
\end{pycode}
\end{region}

\begin{region}{329.3}
\emo{1F527}\textbf{In} \code{wall.py}
\end{region}

\begin{region}{345.1}
In the method collision(), a brick that was hit used to be simply deleted,
e.g., like this:
\end{region}

\begin{region}[pitch=13.2]{364.0}
\begin{pycode}
    \kw{if} Wall.brick[i2, j2]:
        Wall.brick[i2, j2] = 0
\end{pycode}
\end{region}

\begin{region}{394.2}
Now replace this with:
\end{region}

\begin{region}[]{411.5}
\begin{pycode}
    \kw{if} Wall.brick[i2, j2]:
        Wall.brick[i2, j2] = BrickType.effect(Wall.brick[i2, j2])
\end{pycode}
\end{region}

\begin{region}{440.7}
This way, each brick decides for itself what happens to it when it is hit.
\end{region}

\begin{region}{481.2}
\subheading{Task}
\end{region}

\begin{region}{507.9}
\begin{tasks}
\item Test whether bricks of type \code{TWICE} now only disappear on the
  second hit.
\end{tasks}
\end{region}

\begin{region}[x=74.7,right=0,pitch=14.3]{535.3}
\textbf{Pro tip}: To test this specifically, you can \textbf{temporarily set
all bricks to} \code{TWICE} in the \code{init()} method of \code{Wall},
e.g.:
\end{region}

\begin{region}[x=74.7]{568.7}
\begin{pycode}
    Wall.brick[i, j] = \num{4} \cmt{\# ty}
\end{pycode}
\end{region}

\begin{region}[x=74.7,right=0,pitch=13.2]{587.4}
This saves you the tedious search for the right brick -- and you can see
immediately whether the behavior is correct.
\end{region}

\begin{region}[pitch=13.9]{614.7}
\begin{tasks}[start=2]
\item \textbf{The easy case? Think again.}\\
  Actually, everything would now be quite easy: when a brick of type
  \code{BALL} is hit, \textbf{a new ball should come into play}.
\end{tasks}
\end{region}

\begin{region}[x=74.7,right=0,pitch=14.3]{662.2}
\textbf{Sounds simple}: just create a new \code{Ball} object, append it to
\code{Game.balls}, done! \textbf{But think again}. Python unfortunately makes
the \textbf{actually easy case} complicated.
\end{region}

\begin{region}[x=74.7]{699.7}
\emo{1F9EA}\textbf{Experiment: what goes wrong?}
\end{region}

\begin{region}[x=74.7]{721.2}
In \code{bricktype.py}, add the following line at the top:
\end{region}

\begin{region}[x=74.7]{740.2}
\begin{pycode}
\imp{from} ball \imp{import} Ball
\end{pycode}
\end{region}

\begin{region}[x=74.7,right=0]{758.6}
Make no other changes -- and start the game. It will not work.
\end{region}

\nextpage
%% ===================================================================== 8 ==
\begin{region}{69.4}
\heading{Step 6: Cyclic dependencies}
\end{region}

\begin{region}{88.1}
What \textbf{went wrong} at the end of the last step? You have created the
following \textbf{import chain}:
\end{region}

\begin{region}[pitch=13.5]{101.0}
\begin{dashlist}
\item \code{bricktype.py} imports \code{ball.py}
\item \code{ball.py} imports \code{wall.py}
\item \code{wall.py} imports \code{bricktype.py}
\end{dashlist}
\end{region}

\begin{region}{146.3}
\catemo\textbf{And so the cat is chasing its own tail.}
\end{region}

\begin{region}[pitch=14.27]{168.5}
In an ``ideal'' object-oriented language, every class can \textbf{refer to
other classes arbitrarily} -- after all, that is exactly the point of modular
development. But in \textbf{Python}, \textbf{all of the code is executed
immediately} on import -- and because Python does not cleanly separate
\textbf{interface} and \textbf{implementation}, \textbf{cyclic dependencies}
lead to technical problems.
\end{region}

\begin{region}{234.5}
\emo{1FA79}\textbf{Makeshift solutions (not recommended)}
\end{region}

\begin{region}[pitch=13.9]{251.2}
An unclean solution would be to replace (almost) \textbf{all}
\code{from ... import ...} statements\\
with plain \code{import ...} statements and to use, e.g.,
\code{module.Class} later on.\\
This \textbf{circumvents} the problem -- but it is not pretty and breaks with
the style.
\end{region}

\begin{region}[pitch=13.28]{298.0}
Another (even more questionable) option would be to write the
\code{from ball import Ball} directly into the method where the ball is
needed -- e.g., into \code{BrickType.effect}.\\
That would be a quick solution, but honestly: it breaks my heart. On every
call, Python would check whether the module needs to be loaded. Even if this
is technically cached: there is no need to trample all over performance.
\end{region}

\begin{region}{374.2}
\emo{1F9FC}\textbf{A somewhat cleaner solution: controlling the imports}
\end{region}

\begin{region}[pitch=13.5]{395.7}
We break the import cycle by no longer importing \code{BrickType} directly
in \code{wall.py}, but passing it to \code{Wall} at runtime instead.
\end{region}

\begin{region}{432.3}
\emo{1F9F1}\textbf{Changes in} \code{wall.py}
\end{region}

\begin{region}{453.8}
At the top, remove the line: \code{from bricktype import BrickType}
\end{region}

\begin{region}{472.5}
Replace the beginning of the class with:
\end{region}

\begin{region}[pitch=13.2]{491.4}
\begin{pycode}
\kw{class} \fn{Wall}:
    brick = numpy.zeros((Game.brick\_rows, Game.brick\_cols))
\end{pycode}
\end{region}
\begin{region}[]{516.6}
\begin{pycode}
    BrickType = \fn{None}
\end{pycode}
\end{region}

\begin{region}{534.1}
Adapt the method \code{init()} so that \code{BrickType} is passed and
stored:
\end{region}

\begin{region}[pitch=13.2]{553.1}
\begin{pycode}
\kw{def} \fn{init}(BrickType):
    Wall.BrickType = BrickType
\end{pycode}
\end{region}
\begin{region}[]{578.3}
\begin{pycode}
    \kw{for} i \kw{in} \fn{range}(\num{4}, \num{8}):
        \kw{for} j \kw{in} \fn{range}(\num{0}, Game.brick\_cols):
            \kw{if} random.randint(\num{1}, \num{100}) <= \num{20}:
                ty = random.randint(\num{2}, \fn{max}(Wall.BrickType))
            \kw{else}:
                ty = \num{1}
            Wall.brick[i, j] = ty
\end{pycode}
\end{region}

\begin{region}[pitch=13.3]{667.8}
Change the remaining code in \code{wall.py} so that \code{Wall.BrickType} is
always used instead of \code{BrickType}.
\end{region}

\begin{region}{704.1}
\emo{1F527}\textbf{Changes in} \code{breakout.py}
\end{region}

\begin{region}[pitch=13.5]{725.6}
As usual, import \code{from bricktype import BrickType} at the top.\\
Instead of \code{Wall.init()}, now call \code{Wall.init(BrickType)}.
\end{region}

\begin{region}[pitch=14.2]{758.8}
{\bfseries\color{aslabblue}Make the changes as described and test whether
everything still works. If so, you now have a clean, cycle-free design and
are ready for the next step\par}
\end{region}

\nextpage
%% ===================================================================== 9 ==
\begin{region}{69.4}
\heading{Step 7: Back to features}
\end{region}

\begin{region}[pitch=13.4]{88.3}
Now that we have elegantly solved the problem of cyclic dependencies,\\
we can finally import everything we need in \code{bricktype.py} --
\end{region}

\begin{region}{119.3}
(at least everything except from \code{breakout.py} \emoji{1F609}).
\end{region}

\begin{region}{140.5}
For the next step, this is enough:
\end{region}

\begin{region}[pitch=13.2]{159.4}
\begin{pycode}
\imp{from} ball \imp{import} Ball
\imp{from} game \imp{import} Game
\end{pycode}
\end{region}

\begin{region}{198.6}
\emo{1F527}\textbf{Adapting the} \code{effect()} \textbf{method}
\end{region}

\begin{region}{220.1}
Now you can extend the method \code{effect()} in the class \code{BrickType}
as planned:
\end{region}

\begin{region}[pitch=13.2]{239.1}
\begin{pycode}
\kw{def} \fn{effect}(self):
    \kw{if} self == BrickType.TWICE:
\end{pycode}
\end{region}
\begin{region}[]{264.3}
\begin{pycode}
        \kw{return} BrickType.NORMAL
    \kw{elif} self == BrickType.BALL:
        Game.balls.append(Ball(\num{100}, \num{0}, -\num{3}, \num{4}))
        \kw{return} BrickType.NONE
    \kw{else}:
       \kw{return} BrickType.NONE
\end{pycode}
\end{region}

\begin{region}{345.9}
\emo{1F9EA}\textbf{Test}
\end{region}

\begin{region}[pitch=13.4]{362.6}
With this change, destroying a brick of type \code{BALL} should bring a
\textbf{new ball} into the game.
\end{region}

\begin{region}[pitch=13.2]{394.9}
To test this specifically, you can -- as in earlier tasks -- temporarily
initialize the bricks in \code{wall.py} like this:
\end{region}

\begin{region}[]{427.0}
\begin{pycode}
Wall.brick[i, j] = \num{2} \cmt{\# ty}
\end{pycode}
\end{region}

\begin{region}{445.5}
Then you will see immediately whether a new ball is created on a hit.
\end{region}

\begin{region}{472.8}
\subheading{Task}
\end{region}

\begin{region}[pitch=13.9]{492.7}
While testing, you may be pleased only briefly that a new ball comes into
play. Unfortunately, it is created \textbf{above the brick wall} -- and then
makes for a rather boring game, because it simply clears \textbf{bricks from
top to bottom} until a big hole has been punched through.
\end{region}

\begin{region}{557.4}
\emo{1F9E0}\textbf{Idea}
\end{region}

\begin{region}[pitch=13.2]{578.6}
We change the game behavior so that a newly created ball is initially
inactive -- i.e., it flies through the bricks until it is hit by the paddle.
Only then does it become active and can destroy bricks.
\end{region}

\begin{region}{635.7}
\emo{1F527}\textbf{Implementation in} \code{ball.py}
\end{region}

\begin{region}{657.2}
\begin{tasks}
\item In the constructor (\code{__init__}), the ball gets a new attribute:
\end{tasks}
\end{region}

\begin{region}[x=74.7]{675.1}
\begin{pycode}
self.active = \fn{False}
\end{pycode}
\end{region}

\begin{region}[pitch=13.6]{692.7}
\begin{tasks}[start=2]
\item In the \code{move()} method:
  \begin{bullets}
  \item Only call \code{Wall.collision(...)} if \code{self.active == True}
  \item When the ball hits the paddle, set \code{self.active = True}
  \end{bullets}
\end{tasks}
\end{region}

\begin{region}[pitch=14.2]{747.1}
{\bfseries\color{aslabblue}Start the game and check whether a new ball now
first falls through the brick wall\\
and can destroy bricks only after contact with the paddle.\par}
\end{region}

\nextpage
%% ==================================================================== 10 ==
\begin{region}{69.4}
\heading{Step 8: Protective shield}
\end{region}

\begin{region}[pitch=13.2]{95.9}
Now the paddle gets yet another feature: a protective shield that prevents
bomb damage for a certain time.
\end{region}

\begin{region}{139.8}
\emo{1F9F1}\textbf{Preparation in} \code{paddle.py}
\end{region}

\begin{region}[pitch=13.3]{161.3}
First, the class Paddle is extended with a new attribute \code{shield},
similar to \code{damage}.\\
It counts down as long as the shield is active:
\end{region}

\begin{region}[pitch=13.2]{193.5}
\begin{pycode}
\kw{class} \fn{Paddle}:
    \cmt{\# ...}
\end{pycode}
\end{region}
\begin{region}[]{218.7}
\begin{pycode}
    damage = \num{0}
    shield = \num{0}
\end{pycode}
\end{region}

\begin{region}{252.3}
\emo{1F4A5}\textbf{Adapting the method} \code{hit()}
\end{region}

\begin{region}{273.8}
The shield protects against damage -- so \code{hit()} has to be adapted like
this:
\end{region}

\begin{region}[pitch=13.2]{292.8}
\begin{pycode}
\kw{def} \fn{hit}():
    \kw{if} Paddle.damage == \num{0} \kw{and} Paddle.shield == \num{0}:
\end{pycode}
\end{region}
\begin{region}[]{318.0}
\begin{pycode}
        Paddle.damage = \num{90}
        \kw{return} \fn{True}
    \kw{return} \fn{False}
\end{pycode}
\end{region}

\begin{region}{363.6}
\emo{2728}\textbf{New method} \code{set_shield()}
\end{region}

\begin{region}{379.3}
This method activates the shield:
\end{region}

\begin{region}[pitch=13.2]{398.2}
\begin{pycode}
\kw{def} \fn{set\_shield}():
    Paddle.shield = \num{300}
\end{pycode}
\end{region}

\begin{region}{433.0}
\emo{1F3A8}\textbf{Cosmetics in} \code{draw()}
\end{region}

\begin{region}{448.8}
When the shield is active, the paddle is drawn in white, for example:
\end{region}

\begin{region}[pitch=13.2]{467.7}
\begin{pycode}
\kw{def} \fn{draw}():
    \kw{if} Paddle.damage > \num{0}:
\end{pycode}
\end{region}
\begin{region}[]{492.9}
\begin{pycode}
        Paddle.damage -= \num{1}
        \kw{return}
    \kw{if} Paddle.shield > \num{0}:
        Paddle.shield -= \num{1}
        c = (\num{255}, \num{255}, \num{255})
    \kw{else}:
        c  = Paddle.c
    pygame.draw.rect(Game.screen, c,
                    (Paddle.x, Paddle.y, Paddle.w, Paddle.h))
\end{pycode}
\end{region}

\begin{region}{607.1}
\subheading{Task}
\end{region}

\begin{region}{633.8}
\begin{tasks}
\item Make the changes described above in \code{paddle.py}.
\end{tasks}
\end{region}

\begin{region}[pitch=13.3]{647.3}
\begin{tasks}[start=2]
\item Then adapt the method \code{effect()} in \code{bricktype.py} so that
  bricks of type \code{SHIELD} are supported.
\end{tasks}
\end{region}

\begin{region}[x=74.7,right=0,pitch=14.2]{678.1}
You know the procedure by now:
\begin{bullets}
\item \textbf{Think about} what you additionally need to import.
\item \textbf{Think about} how you can quickly test whether the feature
  works.
\end{bullets}
\end{region}

\begin{region}[pitch=17.6]{733.3}
If everything works, the paddle should be invulnerable to bombs for
10~seconds after destroying a \code{SHIELD} brick -- and glow bright white
while doing so \emoji{26A1}
\end{region}

\nextpage
%% ==================================================================== 11 ==
\begin{region}{69.4}
\heading{Step 9: Countermeasures}
\end{region}

\begin{region}[pitch=14.3]{96.9}
\textbf{Wouldn't it be great} if you could also \textbf{actively defend
yourself against bombs} -- that is, simply \textbf{shoot them down} before
they hit you?
\end{region}

\begin{region}{142.0}
\emo{1F9E0}\textbf{Oh dear.}
\end{region}

\begin{region}[pitch=13.7]{158.7}
Yes, that would be great\dots{} But \textbf{should you really, at the last
minute,} build in yet another new feature instead of working through
important items on the to-do list?\\
No. \textbf{You shouldn't.}
\end{region}

\begin{region}{218.1}
\emo{1F937}\textbf{Just like in real life\dots}
\end{region}

\begin{region}[pitch=14.3]{234.8}
\dots{} we do it \textbf{anyway}. And actually, it is not that big a deal --
it should \textbf{work exactly the same way} as shooting a space invader out
of the sky.
\end{region}

\begin{region}{279.9}
\emo{1F6E0}\textbf{The existing code in} \code{breakout.py}
\end{region}

\begin{region}{295.7}
This is how space invaders are currently hit by balls or rockets:
\end{region}

\begin{region}[]{322.3}
\begin{pycode}
\kw{for} s \kw{in} Game.ships:
    \kw{for} item  \kw{in} Game.balls + Game.rockets:
        dist = math.hypot(item.x - s.x, item.y - s.y)
        \kw{if} dist <= \num{15}:
           s.hit = \fn{True}
           item.hit = \fn{True}
\end{pycode}
\end{region}

\begin{region}[pitch=14.2]{407.2}
Thanks to the way of thinking sharpened in the HM (higher mathematics)
lectures, we immediately recognize: if bombs are to be destroyed
\textbf{exactly like} ships, then we simply have to \textbf{iterate over the
bombs} as well.
\end{region}

\begin{region}{462.1}
So it is enough to replace the loop:
\end{region}

\begin{region}[]{488.7}
\begin{pycode}
\kw{for} s \kw{in} Game.ships:
\end{pycode}
\end{region}

\begin{region}{515.1}
with:
\end{region}

\begin{region}[]{541.7}
\begin{pycode}
\kw{for} s \kw{in} Game.ships + Game.bombs:
\end{pycode}
\end{region}

\begin{region}{569.0}
\subheading{Task}
\end{region}

\begin{region}{599.9}
Obvious, right? \textbf{Just do it}. \emoji{1F604}
\end{region}

\begin{region}{629.8}
\heading{Step 10: Because there is still a little room here}
\end{region}

\begin{region}[pitch=13.25]{649.5}
When all space invaders have been shot down, it gets boring -- after all, we
want to play Breakout \textbf{and} Space Invaders. So let us quickly add the
feature that a new space invader appears for every brick that is hit.
\end{region}

\begin{region}{694.9}
In \code{bricktype.py}, add the following line at the very beginning of the
method \code{effect()}:
\end{region}

\begin{region}[]{714.0}
\begin{pycode}
Game.ships.append(Ship(\num{0}))
\end{pycode}
\end{region}

\begin{region}{731.6}
That should actually be enough.\\
Just test it.
\end{region}

\begin{region}{759.9}
If some import thingy is missing, an error message will tell you
\emoji{1F609}
\end{region}

\nextpage
%% ==================================================================== 12 ==
\begin{region}{69.4}
\heading{Step 11: Take a look at the clock \dots}
\end{region}

\begin{region}[pitch=13.7]{87.2}
\dots{} maybe we should start with one of the important features. So let us
pause briefly and think: \textbf{What absolutely belongs in a game like
Breakout + Space Invaders?} Perhaps that you only have a limited number of
lives -- and fight to get to the next level before all lives are used up. In
other words, exactly what makes university studies, with their limited number
of permitted exam attempts, so exciting.
\end{region}

\begin{region}{160.5}
Okay, how do we start? In any case, we need to \textbf{keep track} of two
important things:
\end{region}

\begin{region}[pitch=14.2]{173.9}
\begin{dashlist}
\item the number of \textbf{lives}
\item the current \textbf{level}
\end{dashlist}
\end{region}

\begin{region}{211.1}
\emo{1F9F1}\textbf{In game.py}
\end{region}

\begin{region}{226.9}
Add two new attributes to the class Game:
\end{region}

\begin{region}[pitch=13.2]{244.7}
\begin{pycode}
\kw{class} \fn{Game}:
    \cmt{\# ... other attributes as before}
\end{pycode}
\end{region}
\begin{region}[]{269.9}
\begin{pycode}
    lifes = \num{0}
    level = \num{0}
\end{pycode}
\end{region}

\begin{region}{303.5}
\emo{1F5BC}\textbf{And then: showing where things stand}
\end{region}

\begin{region}[pitch=13.35]{319.5}
So that you can also see what situation you are in, we display \code{level}
and \code{lifes} in the top left corner of the window. Luckily, you have
already seen today how to render text -- so you can, e.g., add the following
at the end of the \code{draw()} method:
\end{region}

\begin{region}[pitch=13.2]{365.2}
\begin{pycode}
    \kw{def} \fn{draw}():
        \cmt{\# ... everything as before}
\end{pycode}
\end{region}
\begin{region}[]{390.4}
\begin{pycode}
        label = f\num{"Level \{Game.level\} Lifes \{Game.lifes\}"}
        text = Game.font.render(label, \fn{True}, (\num{255}, \num{255}, \num{255}))
        text\_rect = text.get\_rect(topleft=(\num{10}, \num{10}))
        Game.screen.blit(text, text\_rect)
\end{pycode}
\end{region}

\begin{region}{448.0}
\emo{1F9EA}\textbf{Test}
\end{region}

\begin{region}[pitch=14.3]{470.2}
Once you have made these two changes, you should see in the top left corner
that you have \textbf{0 lives} and are at \textbf{level 0}. Sounds sobering
at first -- but: \textbf{it is a start}.
\end{region}

\begin{region}{507.6}
\emo{1F501}\textbf{And now: an endless loop}
\end{region}

\begin{region}[pitch=13.4]{528.7}
Now we want to let the game \textbf{keep running robustly}, i.e., build in a
simple restart mechanism:
\end{region}

\begin{region}[pitch=13.33]{554.7}
\begin{dashlist}
\item When you have no lives left, the game starts over from the beginning.
  For example, with 3 lives at level~1
\item When there is no ball left in the game, you lose a life (if you still
  have one), and a new ball comes into play.
\end{dashlist}
\end{region}

\begin{region}{616.6}
\emo{1F527}\textbf{Implementation in} \code{breakout.py}
\end{region}

\begin{region}[pitch=13.5]{637.1}
Start with an empty ball list, \code{Game.balls = []}, and implement the
idea inside the main loop (e.g., directly after \code{while running:}) as
follows:
\end{region}

\begin{region}[]{677.2}
\begin{pycode}
    \kw{if} \fn{len}(Game.balls) == \num{0}:
        \kw{if} Game.lifes > \num{1}:
            Game.lifes -= \num{1}
        \kw{else}:
            Game.lifes = \num{3}
            Game.level = \num{1}
            Wall.init(BrickType)
        Game.balls.append(Ball(\num{400}, \num{0}, -\num{2}, \num{6}))
\end{pycode}
\end{region}

\nextpage
%% ==================================================================== 13 ==
\begin{region}{82.6}
\subheading{Task}
\end{region}

\begin{region}[pitch=13.15]{109.1}
\begin{tasks}
\item[1.] Implement the changes described above and test the game. You
  should now start with 3 lives at level 1 (that is already an
  improvement). When all balls have been lost, you lose a life.
\end{tasks}
\end{region}

\begin{region}[pitch=13.17]{148.6}
\begin{tasks}[start=2]
\item[2.] Great, you can now lose lives -- but not gain any yet. That may be
  how it is in real life, but in a game you want more. For example: that you
  move up to the next level when all bricks have been cleared -- and then
  face an even greater challenge, such as an additional row of bricks.
\end{tasks}
\end{region}

\begin{region}[x=74.7,right=0,pitch=13.3]{207.0}
To do this, change the \code{init} method in \code{wall.py} so that the
number of brick rows corresponds to the current level:
\end{region}

\begin{region}[x=74.7,pitch=13.2]{239.2}
\begin{pycode}
\kw{def} \fn{init}(BrickType):
     Wall.BrickType = BrickType
     Wall.brick = numpy.zeros((Game.brick\_rows, Game.brick\_cols))
     \kw{for} i \kw{in} \fn{range}(\num{4}, \num{4} + Game.level):
         \kw{for} j \kw{in} \fn{range}(\num{0}, Game.brick\_cols):
             \kw{if} random.randint(\num{1}, \num{100}) <= \num{20}:
                 ty = random.randint(\num{2}, \fn{max}(Wall.BrickType))
             \kw{else}:
                 ty = \num{1}
             Wall.brick[i, j] = ty
\end{pycode}
\end{region}

\begin{region}[x=74.7,right=0,pitch=13.2]{377.0}
In \code{breakout.py}, first remove the old initialization with
\code{Wall.init(BrickType)}. Instead, you can now regularly check in the main
loop how many bricks are still in play. When none are left, you move up a
level:
\end{region}

\begin{region}[x=74.7,pitch=13.2]{422.3}
\begin{pycode}
\kw{if} Wall.brick.sum() == \num{0}:
    Game.level += \num{1}
    Wall.init(BrickType)
\end{pycode}
\end{region}

\begin{region}[x=74.7,right=0,pitch=13.2]{467.7}
The method \code{sum} comes from \code{numpy}, which -- as you know -- is a
wrapper for a very efficient C/C++ library for numerical linear algebra.
Computing the sum of all matrix elements is extremely efficient. So it is
not worth counting separately on the Python side how many bricks are still
left. The code is short, clear, and performant.
\end{region}

\begin{region}[x=74.7,right=0,pitch=13.2]{524.9}
Make these changes and test the game. You now start with one row of bricks.
If you manage to clear it completely, you move up a level and face a new
wall with an additional row. If you have lost all lives, you start again
with 3 lives at level 1.
\end{region}

\begin{region}{591.7}
\heading{Step 12: Since there is still room, here's something for your ears}
\end{region}

\begin{region}[pitch=13.2]{609.4}
Unfortunately, the PCs in the lab rooms have no speakers. But there are
headphones! And anyone who feels like giving the game a bit more punch can
now add sound effects.
\end{region}

\begin{region}[pitch=13.2]{640.1}
With pygame, this is very easy. A sound object is created from an audio file
-- and you can play it with .play() whenever you feel like it.
\end{region}

\begin{region}{670.8}
For example, in game.py in the init method:
\end{region}

\begin{region}[pitch=13.2]{688.7}
\begin{pycode}
Game.exploding = pygame.mixer.Sound(\num{"data/exploding.wav"})
Game.bounce = pygame.mixer.Sound(\num{"data/bounce.mp3"})
\end{pycode}
\end{region}
\begin{region}[]{713.9}
\begin{pycode}
Game.launch = pygame.mixer.Sound(\num{"data/launch.wav"})
\end{pycode}
\end{region}

\begin{region}{730.0}
When, e.g., a rocket is fired, the matching clip can then be played like
this:
\end{region}

\begin{region}[]{747.8}
\begin{pycode}
Game.launch.play()
\end{pycode}
\end{region}

\begin{region}[pitch=13.1]{765.5}
There is no more room for explanations here -- but that does not matter,
because at this point, your enthusiasm for experimenting is called for. Think
about where else sound effects could be played.
\end{region}

\nextpage
%% ==================================================================== 14 ==
\begin{region}{69.4}
\heading{Step 13: Oh yes, the repository \dots}
\end{region}

\begin{region}[pitch=17.5]{87.2}
If anyone has made it this far: it is about time to commit the
changes to the repository -- unless you have already come up with that idea
yourselves anyway \emoji{1F609}
\end{region}

\begin{region}[pitch=13.2]{133.9}
\begin{pycode}
git add *.py
git commit -a -m "about time"
git push
\end{pycode}
\end{region}

\begin{region}{200.5}
\heading{Step 14: And then there is the matter of the controller}
\end{region}

\begin{region}[pitch=13.15]{218.2}
In the lab, we have provided you (unless we forgot) with a game controller
with a BNO055 sensor. With it, you can steer the paddle casually with a
flick of your wrist.
\end{region}

\begin{region}[pitch=13.18]{270.9}
It should be surprising -- or maybe not -- that you did not have to adapt
your Python code for reading the serial interface at all. Why? Because it
makes no difference at all whether the microcontroller reads a potentiometer
or another sensor -- as long as it sends the data over the serial interface
using the same protocol, everything works as before.\\
If you want the matching microcontroller code for the BNO055:
\end{region}

\begin{region}[]{350.2}
\begin{pycode}
git clone git@github.com:michael-lehn/asp-ulm-bno055.git
\end{pycode}
\end{region}

\begin{region}{376.6}
The Makefile works the way you know it by now:
\end{region}

\begin{region}[pitch=13.5]{402.5}
\begin{dashlist}
\item \code{make} compiles the code
\item \code{make upload} compiles and flashes
\item \code{make monitor} starts the serial monitor
\end{dashlist}
\end{region}

\begin{region}{470.7}
\heading{Step 15: And here there would still be room \dots}
\end{region}

\begin{region}[pitch=13.2]{493.9}
\dots{} for a long to-do list with refactoring, clean code, a menu, a high
score, and options.\\
Or for a few closing words about the path we have taken.
\end{region}

\begin{region}{524.6}
But honestly: I like neither to-do lists nor goodbyes \emoji{1F604}
\end{region}

\end{document}
